\chapter{Conclusion and Future Work}

\section{Summary}

This report has documented the full software-engineering lifecycle of the
\textbf{Smart Home Energy Optimizer (SHEO)}, an AI-assisted monitoring and
control platform designed to help residential occupants reduce electricity
consumption and visualise the financial impact of their energy-saving decisions
in Vietnamese Dong (VND).

SHEO is delivered as a three-tier web application: a FastAPI back-end organised
into layered services and a capability-driven device abstraction, a React
single-page front-end, and a mock device simulator that enables end-to-end
demonstration without physical hardware.  The platform's principal capabilities
are:

\begin{itemize}[noitemsep]
  \item \textbf{Capability-driven monitoring and control.} Every device exposes a
        capability schema that drives both the API and the front-end widget
        automatically, so adding a new device type requires no UI code changes.
  \item \textbf{Explainable WHEN--THEN automation rules.} Residents author rules
        for their own unit in structured natural language; the rule
        engine interprets and audits them rather than executing opaque scripts.
  \item \textbf{Habit recommendations.} A recommendation engine identifies
        high-consumption patterns and proposes WHEN--THEN rules; residents review
        and accept each suggestion individually.
  \item \textbf{VND bill-saving estimation.} Before accepting a recommendation,
        the user sees an estimated monthly saving in VND.  Accepted rules
        accumulate in a cycle-savings dashboard, closing the feedback loop between
        behaviour change and financial outcome.
\end{itemize}

The system is fully implemented and verified: \textbf{61 automated tests} (unit,
integration, system, and acceptance levels) pass against the current codebase,
and all functional requirement families listed in the SRS are represented in the
test suite's traceability matrix.

\section{Requirements Satisfaction}

The client's single highest-priority acceptance criterion was:

\begin{quote}
\textit{The Customer can see the actual VND amount saved by a suggested
automation rule before deciding whether to accept it.}
\end{quote}

This criterion is satisfied through two complementary mechanisms.  First, the
\emph{estimate-before-accept} flow computes a prospective monthly saving (in
VND) at the point of recommendation review; the user sees the figure before
committing.  Second, the \emph{cycle-savings dashboard} aggregates realised
savings for each accepted rule over the current billing cycle, providing
post-hoc validation that the estimate was meaningful.

Table~\ref{tab:req-satisfaction} maps each functional requirement family to its
implementation.  Collectively, these families are exercised by the suite of
61 automated tests, whose per-requirement traceability is detailed in the
appendix.

\begin{table}[htbp]
  \centering
  \caption{Functional requirement families and their satisfaction status}
  \label{tab:req-satisfaction}
  \begin{tabularx}{\textwidth}{@{} l X X @{}}
    \toprule
    \textbf{Req.\ family} & \textbf{Description} & \textbf{Key component} \\
    \midrule
    REQ-4.1 & Real-time device monitoring & The device-listing and monitoring service with its capability widget \\
    REQ-4.2 & Remote device control & The device-control operation \\
    REQ-4.3 & WHEN--THEN automation rules & The rule engine \\
    REQ-4.4 & Habit recommendations & The recommendation service and its savings estimate \\
    REQ-4.5 & VND bill-saving dashboard & Cycle-savings aggregation in the bill-saving (optimisation) service \\
    \bottomrule
  \end{tabularx}
\end{table}

Non-functional requirements were addressed as follows: response-time
constraints (NFR-PER) are enforced by asynchronous request handlers in the
application server and in-memory caching; safety constraints (NFR-SAF) prevent
the system from automatically powering off safety-critical devices; security
constraints (NFR-SEC) are upheld by role-based access guards (the
administrator-authorisation guards) and the IDOR vulnerability discovered during
the adversarial review was patched and regression-tested.

\section{Course-Concept Coverage}

Table~\ref{tab:course-coverage} provides a candid assessment of how each
lecture topic in IT3180E is reflected in the SHEO project artifacts.

{\small
\begin{xltabular}{\textwidth}{@{} l l X l @{}}
  \caption{IT3180E lecture-topic coverage by SHEO project}\label{tab:course-coverage}\\
    \toprule
    \textbf{Lecture} & \textbf{Topic} & \textbf{Where addressed in this report} & \textbf{Coverage} \\
    \midrule
    \endfirsthead
    \toprule
    \textbf{Lecture} & \textbf{Topic} & \textbf{Where addressed in this report} & \textbf{Coverage} \\
    \midrule
    \endhead
    C01 & SE overview; stakeholder roles
        & Chapter~1: Client / Customer / Resident model, a general framing of the software-engineering lifecycle, quality focus, Function--Cost--Time triangle
        & Strong \\[4pt]
    C02 & Software development process
        & Chapter~2: incremental model with agile-style iterations; ISO/IEC~9126 quality attributes mapped to NFRs
        & Partial \\[4pt]
    C03 & Agile / Scrum
        & Chapter~2: working-software discipline maintained throughout; no formal Scrum roles or sprint artifacts documented
        & Light \\[4pt]
    C04 & Software project management
        & Chapter~3: Work Breakdown Structure, effort-estimation table, CPM task network with critical path, Gantt schedule, and a risk register scored by Probability~$\times$~Impact with Avoid/Transfer/Mitigate/Accept responses
        & Strong \\[4pt]
    C05 & Software configuration management
        & Chapter~8: dedicated CM plan identifying CIs, baselines, and Git workflow
        & Partial \\[4pt]
    C06 & Requirements analysis and UML
        & Chapter~4: SRS, FR/NFR taxonomy, Use-Case Diagram, ERD, state-machine diagrams, activity-style user flow
        & Strong \\[4pt]
    C07 & Software design and UI/UX
        & Chapter~5 (architecture, detailed design, and patterns) and Chapter~6 (UI/UX): context-of-use, user profiles, sitemap, user flow, wireframes, usability evaluation
        & Strong \\[4pt]
    C08 & Software construction
        & Chapter~7: coding-style conventions, adversarial code review as formal inspection analogue, refactoring findings
        & Partial \\[4pt]
    C09 & Quality assurance: testing and maintenance
        & Chapter~9: V-model mapping, black-box equivalence partitioning, boundary-value analysis, decision-table testing, white-box CFG examples, regression suite; maintenance classification
        & Partial \\
    \bottomrule
  \end{xltabular}
}

The project is strongest where the course emphasises artefact production.
\textbf{C01} coverage is strong because the stakeholder model is rendered
end-to-end from the glossary through to the role enumeration in the
implementation.
\textbf{C04} coverage is strong because Chapter~3 produced the full
project-management suite --- Work Breakdown Structure, effort estimation, a CPM
task network with an identified critical path, a Gantt schedule, and a risk
register scored by Probability~$\times$~Impact with explicit
Avoid/Transfer/Mitigate/Accept responses.
\textbf{C06} coverage is strong because the project produced a full SRS together
with a complementary suite of UML diagrams (use-case, ERD, state machines).
\textbf{C07} coverage is strong across both the architectural-design half and the
UI/UX half.  Coverage is partial for \textbf{C02} (the incremental process model
was followed in practice but not formally justified in a process document),
\textbf{C05} (a dedicated CM plan exists but IEEE~828 compliance and the exact
branching strategy are not fully cross-referenced from the main body),
\textbf{C08} (the adversarial review activity occurred and its findings were
corrected, but a formal inspection report was not produced as a standalone
artefact), and \textbf{C09} (white-box CFG examples were added but a complete
branch-coverage report was out of scope).  Coverage is lightest for \textbf{C03},
which concerns the team-level Scrum process that is inherently difficult to
demonstrate through a single design-and-build deliverable.

\section{Lessons Learned}

\begin{itemize}[leftmargin=1.5em, itemsep=4pt]

  \item \textbf{A single source of truth yields cumulative benefit.}
        Anchoring every API endpoint, front-end widget, and test assertion to the
        capability schema eliminated an entire class of synchronisation bugs.
        When the schema changed, the ripple effect was immediately visible
        rather than silently divergent.

  \item \textbf{Separating requirements from design is not bureaucracy.}
        Maintaining a discrete SRS before writing the SDD forced the team to
        confront ambiguities during elicitation rather than during coding.  The
        \emph{what} document also served as the primary input for the
        traceability matrix, making acceptance-test authorship straightforward.

  \item \textbf{An adversarial review surfaces defects that a cooperative walkthrough does not.}
        Running a structured correctness and security review from a
        non-author perspective (analogous to a formal inspection) surfaced a
        cross-unit IDOR vulnerability and an unvalidated rule-update path that
        had survived all prior testing.  Both findings were patched and
        regression-tested before the final submission.

  \item \textbf{Client elicitation must be documented, not recalled.}
        The highest-priority acceptance criterion --- VND savings visibility ---
        was stated clearly in the initial elicitation session; because it was
        captured in the SRS, the team could verify against it mechanically.
        Requirements that were held only in memory tended to drift or be
        deprioritised under time pressure.

  \item \textbf{The process model should be named early, not inferred later.}
        SHEO was built incrementally --- a working vertical slice first, then
        features fanned out --- which maps naturally to the Incremental model.
        Not naming this decision explicitly in the process documentation made the
        C02 coverage weaker than the actual practice warranted.

  \item \textbf{Mock-before-hardware is a legitimate and traceable scope decision.}
        Decoupling the device-adapter interface from physical hardware
        from day one allowed the entire system to be demonstrated and tested
        without IoT equipment.  The decision was justified in the feasibility
        study and respected throughout; it did not compromise the architectural
        validity of the design.

\end{itemize}

\section{Future Work}

Several directions would extend SHEO from a demonstration system to a
production-grade platform.

\begin{description}[style=nextline, leftmargin=1.5em, itemsep=6pt]

  \item[Real-hardware integration via the device-adapter contract]
        The existing device-adapter interface and capability schema were
        designed with physical hardware in mind.  Replacing the mock
        implementations with drivers for standard IoT messaging protocols
        would require no changes to the API or front-end layers.  A
        priority integration target is the smart air-conditioner controller,
        which accounts for the majority of residential electricity consumption in
        the Vietnamese context.

  \item[Broader device and capability coverage]
        Current device types cover air conditioners, lighting, and socket plugs.
        Extending the capability registry to include water heaters, inverter
        washing machines, and EV chargers would increase the platform's
        addressable saving potential without requiring architectural changes,
        since all device behaviour is expressed through the existing capability
        schema.

  \item[Horizontal scalability]
        The present deployment runs as a single application-server process.
        Introducing a message broker between device adapters and the rule
        engine, combined with multiple stateless API workers behind a load
        balancer,
        would support multi-building deployments and eliminate the single point
        of failure identified in the NFR-SAF analysis.

  \item[Deeper project-management instrumentation]
        The project-management plan in Chapter~3 already establishes a Work
        Breakdown Structure, an effort-estimation table, a CPM schedule with an
        identified critical path, a Gantt chart, and a risk register.  A
        follow-on phase could deepen these artefacts with earned-value tracking
        and periodic re-estimation as the project progresses.  Adopting a formal
        Scrum process with a maintained Product Backlog, two-week sprints, and a
        Definition of Done would additionally close the C03 gap and improve
        team-level visibility during development.

  \item[White-box coverage gate]
        The current test suite emphasises black-box equivalence partitioning and
        boundary-value analysis.  Adding an automated branch-coverage gate with a
        minimum-coverage threshold would provide a C09-compliant white-box
        verification gate and prevent coverage regression as the codebase grows.

  \item[Richer analytics and personalised insights]
        The cycle-savings dashboard currently aggregates savings by device and
        rule.  Future work could incorporate time-series anomaly detection to
        flag unexpected consumption spikes, comparative benchmarking against
        similar household profiles, and personalised rule suggestions derived
        from appliance-usage clustering --- moving SHEO from reactive reporting
        toward proactive energy intelligence.

\end{description}
